\section{Abstract}

This paper presents the mathematical foundations underlying a two-body Newtonian gravitational orbit simulator implemented in C++. The simulation models a planet orbiting a star in a two-dimensional plane using classical mechanics and a symplectic numerical integration scheme known as the Velocity Verlet algorithm. All relevant physical principles and formulae are derived and explained in full, from Newton's Law of Universal Gravitation through to the step-by-step integration procedure used to evolve the system in time.

The motivation for this investigation is to demonstrate that even a relatively simple numerical method, when applied correctly, can faithfully reproduce the qualitative behaviour predicted by Newtonian orbital mechanics — including elliptical, circular, parabolic, and hyperbolic trajectories — without the need for an analytic solution at every timestep.

Key topics covered include: the derivation of gravitational acceleration from first principles, the calculation of orbital and escape velocities, the determination of an appropriate initial velocity direction via the perpendicular unit vector, and the energy-conserving properties of the Velocity Verlet integrator. Where appropriate, concrete numerical examples are provided using the specific initial conditions of the simulation to ground the mathematical exposition in the actual implementation.